\documentclass[12pt]{article}
\usepackage[utf8]{inputenc}
\usepackage[T1]{fontenc}
\usepackage{hyperref}
\title{Wymysorys Language Resources for K-12 Curriculum Development}
\author{Compiled by Claude Sonnet 5.5 (claude-sonnet-5-5) for Andr\'as Kornai}
\date{}
\begin{document}
\maketitle

\section{Introduction}
Wymysorys (Vilamovian; ISO 639-3 wym; Glottocode wymy1235) is a West Germanic language of the town of Wilamowice in southern Poland. UNESCO classifies it as critically endangered. Speaker estimates vary from about 20 native speakers (2017, Wikipedia) to at most fifty (University of Warsaw presentation). A revitalisation movement has been active since about 2012, including teaching to local children.

\section{Primary Educational Resources}
\begin{itemize}
\item Polish--Vilamovian phrasebook from the University of Warsaw, with 24 audio files read by Ti\"oma Fum D\"okter (numbers, family, travel and other topics): \url{https://culturalcontinuity.al.uw.edu.pl/resource/polish-vilamovian-phrasebook/}
\item Wikitongues blog post on revitalising Wymysorys and identity in southern Poland: \url{https://wikitongues.org/blog/celebrating-wymysorys-revitalizing-a-language-and-identity-in-southern-poland/}
\end{itemize}

\section{Linguistic and Reference Materials}
\begin{itemize}
\item Wymysorys (Wikipedia): \url{https://en.wikipedia.org/wiki/Wymysorys}
\item Andrason and Kr\'ol, dictionary instalment (letter Z), \emph{Annales Neophilologiarum} 2013: \url{https://bazhum.muzhp.pl/media/texts/annales-neophilologiarum/2013-numer-7/annales_neophilologiarum-r2013-t-n7-s5-29.pdf}
\item Andrason presentation (Duke Slavic centers): \url{https://slaviccenters.duke.edu/sites/slaviccenters.duke.edu/files/media_items_files/andrason.original.pdf}
\item The Language Closet, ``Obscure languages: Wymysorys'': \url{https://thelanguagecloset.com/2021/05/29/obscure-languages-wymysorys/}
\item Languagehat post on Wymysorys: \url{https://languagehat.com/wymysorys/}
\end{itemize}
Search results state that the Wiktionary Wymysorys dictionary grew to over 7,000 words by 2018 as part of the ``Wilamowice'' project, but I saw no direct URL for it.

\section{Religious and Community Resources}
\begin{itemize}
\item Poster ``Ti\"oma Fum D\"okter'' from the Engaged Humanities project, University of Warsaw: \url{https://engagedhumanities.al.uw.edu.pl/wp-content/uploads/2019/03/polaczone-postery.pdf}
\item University of Warsaw cultural continuity project page: \url{https://culturalcontinuity.al.uw.edu.pl/?p=1083}
\end{itemize}

\section{Digital and Audio Resources}
The phrasebook above includes audio. The Wiktionary dictionary is mentioned in the Curriculum section below; no direct URL was seen.

\section{Curriculum Development Context}
Wymysorys has a standardised Latin-based alphabet, a dictionary effort, a phrasebook with audio and reported teaching to children in Wilamowice since about 2012. No complete K--12 curriculum was found. Because few elderly native speakers remain, recording and involving them in materials development is urgent. Textbooks and graded readers would need to be produced with the community and the Warsaw researchers.

\end{document}